% part2.tex — 点粒子的辐射反作用(原 script.tex 第 405-1069 行重构;公式全部不编号,记号全部带量纲)

\subsection{点粒子的辐射反作用}

上一节把电荷轨道当作已知源,完成了固定轨道上的频谱分析。现在要让轨道本身由电磁场决定,就必须回答一个更基本的问题:点电荷的推迟场在世界线上发散,哪一部分应当解释为粒子惯性的重定义,哪一部分才是有限的辐射反作用?这个问题不能从非相对论的“附加质量”概念猜测,而应先在四维时空中完成 Dirac 的场分解与守恒量计算,求得世界线上的有限正则场后,再在有限尺寸粒子的极限下读出熟悉的电磁质量与延迟自作用。

\subsubsection{四维 Green 函数与推迟、超前场}

解决上述问题的第一步,是把推迟场、超前场及其分解在四维时空中写清楚。为此统一使用国际单位制和号差 $+2$,固定记号
\begin{equation*}
 g_{\mu\nu}=\diag(1,-1,-1,-1),
 \qquad
 x^\mu=(ct,\vb x),
 \qquad
 A^\mu=(\Phi/c,\vb A),
 \qquad
 J^\mu=(c\rho,\vb J).
\end{equation*}
粒子世界线记为 $z^\mu(\tau)$,其四速度、四加速度和四维 jerk 分别为
\begin{equation*}
 u^\mu={\d z^\mu\over\d\tau},
 \qquad
 a^\mu={\d u^\mu\over\d\tau},
 \qquad
 \dot a^\mu={\d a^\mu\over\d\tau},
 \qquad
 u^2=c^2,
 \quad
 u\cdot a=0.
\end{equation*}
大写 $J^\mu$ 始终表示电流,点表示对固有时求导;世界管半径用 $\ell$,近世界线的类空位移用 $\xi^\mu$。记号约定小结:$\gamma$ 只保留给三维相对论中的 Lorentz 因子,不再兼任空间位移。

在 Lorenz 规范下,
\begin{equation*}
 F^{\mu\nu}=\partial^\mu A^\nu-\partial^\nu A^\mu,
 \qquad
 \partial_\mu A^\mu=0,
 \qquad
 \Box A^\mu=\mu_0J^\mu,
\end{equation*}
其中
\begin{equation*}
 \Box=\partial_\mu\partial^\mu
 ={\partial^2\over\partial(x^0)^2}-\nabla^2
 ={1\over c^2}{\partial^2\over\partial t^2}-\nabla^2.
\end{equation*}
为避免在一开始把因果边界条件隐藏在三维推迟势中,先求四维 Green 函数。令
\begin{equation*}
 \Box G(x)=\delta^{(4)}(x),
 \qquad
 G(x)=\int{\d^4k\over(2\pi)^4}\,\widetilde G(k)e^{-ik\cdot x}.
\end{equation*}
由于
\begin{equation*}
 \Box e^{-ik\cdot x}=-k^2e^{-ik\cdot x},
 \qquad
 k^2=(k^0)^2-\vb k^2,
\end{equation*}
动量空间 Green 函数必须满足
\begin{equation*}
 -k^2\widetilde G(k)=1.
\end{equation*}
分母在 $k^0=\pm|\vb k|$ 处有极点;不同的极点绕行规定给出不同的边界条件。推迟与超前 Green 函数分别写成
\begin{align*}
 G_{\mathrm{ret}}(x)
 &=-\int{\d^4k\over(2\pi)^4}
 {e^{-ik\cdot x}\over k^2+i0\,k^0},\\
 G_{\mathrm{adv}}(x)
 &=-\int{\d^4k\over(2\pi)^4}
 {e^{-ik\cdot x}\over k^2-i0\,k^0}.
\end{align*}
对 $k^0$ 作围道积分。对推迟解,当 $x^0>0$ 时闭合下半平面,两个极点都被压到下方;当 $x^0<0$ 时闭合上半平面而没有极点。因此
\begin{align*}
 G_{\mathrm{ret}}(x)
 &=\theta(x^0)
 \int{\d^3k\over(2\pi)^3}
 e^{i\vb k\cdot\vb x}
 {\sin(|\vb k|x^0)\over|\vb k|},\\
 G_{\mathrm{adv}}(x)
 &=-\theta(-x^0)
 \int{\d^3k\over(2\pi)^3}
 e^{i\vb k\cdot\vb x}
 {\sin(|\vb k|x^0)\over|\vb k|}.
\end{align*}
令 $r=|\vb x|$,利用
\begin{equation*}
 \int\d\Omega_{\vb k}\,e^{i\vb k\cdot\vb x}
 =4\pi{\sin(kr)\over kr},
\end{equation*}
推迟解的径向积分为
\begin{align*}
 G_{\mathrm{ret}}(x)
 &= {\theta(x^0)\over2\pi^2r}
 \int_0^\infty\d k\,\sin(kr)\sin(kx^0)\\
 &= {\theta(x^0)\over4\pi r}
 \left[\delta(x^0-r)-\delta(x^0+r)\right]\\
 &= {\delta(x^0-r)\over4\pi r}.
\end{align*}
这里第二步用到了 $\int_0^\infty\d k\,\sin(ka)\sin(kb)=(\pi/2)[\delta(a-b)-\delta(a+b)]$,最后一步利用了 $x^0>0$ 与 $r\geq0$。同理
\begin{equation*}
 G_{\mathrm{adv}}(x)={\delta(x^0+r)\over4\pi r}.
\end{equation*}
由于
\begin{equation*}
 \delta(x^2)
 =\delta((x^0)^2-r^2)
 ={\delta(x^0-r)+\delta(x^0+r)\over2r},
\end{equation*}
上面两式又可写成显式协变的光锥形式
\begin{equation*}
 G_{\mathrm{ret/adv}}(x)
 ={1\over2\pi}\theta(\pm x^0)\delta(x^2).
\end{equation*}
因此推迟与超前 Green 函数是同一个波动算符的两个基本解;它们只在对光锥两支的支集选择上不同。

给定任意守恒电流,两个相应的特解为
\begin{equation*}
 A_{\mathrm{ret/adv}}^\mu(x)
 =\mu_0\int G_{\mathrm{ret/adv}}(x-x')J^\mu(x')\d^4x'.
\end{equation*}
它们满足同一个非齐次方程,故半和仍含源而半差无源:
\begin{equation*}
 \Box{A_{\mathrm{ret}}^\mu+A_{\mathrm{adv}}^\mu\over2}
 =\mu_0J^\mu,
 \qquad
 \Box{A_{\mathrm{ret}}^\mu-A_{\mathrm{adv}}^\mu\over2}=0.
\end{equation*}

点电荷沿世界线 $z^\mu(\tau)$ 运动时,四维电流为
\begin{equation*}
 J^\mu(x)=qc\int_{-\infty}^{\infty}
 u^\mu(\tau')\delta^{(4)}(x-z(\tau'))\d\tau'.
\end{equation*}
把它代入上面的势公式,记
\begin{equation*}
 R^\mu(\tau')=x^\mu-z^\mu(\tau'),
 \qquad
 \Phi(\tau')=R^2(\tau'),
\end{equation*}
则
\begin{equation*}
 {\d\Phi\over\d\tau'}
 =2R\cdot{\d R\over\d\tau'}
 =-2R\cdot u.
\end{equation*}
利用单变量 delta 函数公式
\begin{equation*}
 \delta(\Phi(\tau'))
 =\sum_i{\delta(\tau'-\tau_i)\over|\Phi'(\tau_i)|}
 =\sum_i{\delta(\tau'-\tau_i)\over2|R\cdot u|_{\tau_i}},
\end{equation*}
其中 $\tau_i$ 是 $R^2=0$ 的根,便得到
\begin{align*}
 A_{\mathrm{ret}}^\mu(x)
 &= {\mu_0qc\over4\pi}
 \left.{u^\mu\over R\cdot u}\right|_{\tau_{\mathrm{ret}}},\\
 A_{\mathrm{adv}}^\mu(x)
 &=-{\mu_0qc\over4\pi}
 \left.{u^\mu\over R\cdot u}\right|_{\tau_{\mathrm{adv}}}.
\end{align*}
这里推迟根位于过去光锥,满足 $R\cdot u>0$;超前根位于未来光锥,满足 $R\cdot u<0$,故第二式的负号来自绝对值本身,并不是额外的约定。

实际场是同一源方程的某个物理解。任取常数 $\alpha$,它都可写成
\begin{equation*}
 A_{\mathrm{act}}^\mu
 =\alpha A_{\mathrm{ret}}^\mu
 +(1-\alpha)A_{\mathrm{adv}}^\mu
 +A_{\mathrm{hom}}^\mu,
\end{equation*}
其中 $A_{\mathrm{hom}}^\mu$ 是齐次解。对同一个 $A_{\mathrm{act}}^\mu$ 重新分组,定义
\begin{align*}
 A_{\mathrm{in}}^\mu
 &=(1-\alpha)(A_{\mathrm{adv}}^\mu-A_{\mathrm{ret}}^\mu)
 +A_{\mathrm{hom}}^\mu,\\
 A_{\mathrm{out}}^\mu
 &=\alpha(A_{\mathrm{ret}}^\mu-A_{\mathrm{adv}}^\mu)
 +A_{\mathrm{hom}}^\mu,
\end{align*}
则
\begin{equation*}
 A_{\mathrm{act}}^\mu
 =A_{\mathrm{ret}}^\mu+A_{\mathrm{in}}^\mu
 =A_{\mathrm{adv}}^\mu+A_{\mathrm{out}}^\mu,
 \qquad
 A_{\mathrm{out}}^\mu-A_{\mathrm{in}}^\mu
 =A_{\mathrm{ret}}^\mu-A_{\mathrm{adv}}^\mu.
\end{equation*}
因此 $A_{\mathrm{in}}^\mu$ 与 $A_{\mathrm{out}}^\mu$ 都是无源场。对场强作同样的分解,定义 Dirac 奇异场和辐射场
\begin{equation*}
 F_{\mathrm S}^{\mu\nu}
 ={F_{\mathrm{ret}}^{\mu\nu}+F_{\mathrm{adv}}^{\mu\nu}\over2},
 \qquad
 F_{\mathrm R}^{\mu\nu}
 ={F_{\mathrm{ret}}^{\mu\nu}-F_{\mathrm{adv}}^{\mu\nu}\over2},
\end{equation*}
并把在世界线上将保持有限的总正则场记为
\begin{equation*}
 F_{\mathrm{reg}}^{\mu\nu}
 \equiv F_{\mathrm{in}}^{\mu\nu}+F_{\mathrm R}^{\mu\nu}.
\end{equation*}
于是同一个实际场具有连续的等价改写链
\begin{equation*}
 F_{\mathrm{act}}^{\mu\nu}
 =F_{\mathrm{ret}}^{\mu\nu}+F_{\mathrm{in}}^{\mu\nu}
 =F_{\mathrm{adv}}^{\mu\nu}+F_{\mathrm{out}}^{\mu\nu}
 =F_{\mathrm S}^{\mu\nu}+F_{\mathrm{reg}}^{\mu\nu}.
\end{equation*}
后面的近世界线计算将证明:$F_{\mathrm S}^{\mu\nu}$ 精确收集推迟场和超前场共有的局部发散项,而 $F_{\mathrm R}^{\mu\nu}$ 在世界线上有有限极限。超前场在这里只是局部奇异性分解的数学工具;实际散射问题的因果边界条件仍由 $F_{\mathrm{in}}^{\mu\nu}$ 决定。

\subsubsection{近世界线展开：推迟与超前场的共同奇异部分}

上面的等价改写链把问题归结为:读出 $F_{\mathrm S}^{\mu\nu}$ 的局部奇性与 $F_{\mathrm R}^{\mu\nu}$ 的世界线极限,这要求在近世界线区域把推迟场与超前场统一展开。以下全部使用带量纲的四速度、四加速度和 jerk,并以固有时为参数;近线小量是管半径 $\ell$。由 $u^2=c^2$ 连续微分得到
\begin{equation*}
 u^2=c^2,
 \qquad
 u\cdot a=0,
 \qquad
 u\cdot\dot a=-a^2.
\end{equation*}
前两条说明四加速度恒与四速度正交;第三条由 $u\cdot a=0$ 再求导得 $u\cdot\dot a+a^2=0$,是 jerk 与加速度的内积关系,在下面的分母展开中会反复用到。

先把上面的 Liénard--Wiechert 势公式对观测点求导。由光锥条件 $R^2=0$,
\begin{equation*}
 0=\partial_\alpha(R^2)
 =2R_\alpha-2(R\cdot u){\partial\tau'\over\partial x^\alpha},
\end{equation*}
故
\begin{equation*}
 {\partial\tau'\over\partial x^\alpha}
 ={R_\alpha\over R\cdot u}.
\end{equation*}
将上式用于势 $A^\beta=\mu_0qc\,u^\beta/(4\pi\,R\cdot u)$ 的导数,链式法则给出
\begin{equation*}
 \partial_\alpha A^\beta
 ={\mu_0qc\over4\pi}{1\over R\cdot u}
 R_\alpha{\d\over\d\tau'}
 \left({u^\beta\over R\cdot u}\right),
\end{equation*}
再反对称化。利用 $\d R^\alpha/\d\tau'=-u^\alpha$ 把对 $u^\beta/(R\cdot u)$ 的求导转移到 $R^\alpha u^\beta$ 上:
\begin{equation*}
 R^\alpha{\d\over\d\tau'}\left({u^\beta\over R\cdot u}\right)
 ={\d\over\d\tau'}\left({R^\alpha u^\beta\over R\cdot u}\right)
 +{u^\alpha u^\beta\over R\cdot u},
\end{equation*}
组合成反对称后末项的贡献为 $u^{[\alpha}u^{\beta]}=0$,并代回 $\mu_0qc/(4\pi)=q/(4\pi\varepsilon_0c)$,推迟与超前场可统一写成
\begin{equation*}
 F_\eta^{\mu\nu}(x)
 =\left.
 \eta{q\over4\pi\varepsilon_0c}
 {1\over R\cdot u}{\d\over\d\tau'}
 \left({R^{[\mu}u^{\nu]}\over R\cdot u}\right)
 \right|_{\tau'=\tau_\eta},
 \qquad
 \eta=\begin{cases}+1,&\mathrm{ret},\\-1,&\mathrm{adv}.\end{cases}
\end{equation*}
其中 $X^{[\mu}Y^{\nu]}=X^\mu Y^\nu-X^\nu Y^\mu$。

在世界线上固定一点 $\tau_0$。取观测点位于与 $u^\mu(\tau_0)$ 正交的瞬时静止超平面:
\begin{equation*}
 x^\mu=z^\mu(\tau_0)+\xi^\mu,
 \qquad
 \xi\cdot u=0,
 \qquad
 \xi^2=-\ell^2,
 \qquad
 \xi^\mu=\ell n^\mu,
 \quad
 n^2=-1.
\end{equation*}
令源点参数为
\begin{equation*}
 \tau'=\tau_0-\sigma.
\end{equation*}
在 $\tau_0$ 处 Taylor 展开世界线和四速度:
\begin{align*}
 z^\mu(\tau_0-\sigma)
 &=z^\mu(\tau_0)-u^\mu\sigma
 +{1\over2}a^\mu\sigma^2
 -{1\over6}\dot a^\mu\sigma^3
 +O(\sigma^4),\\
 u^\mu(\tau_0-\sigma)
 &=u^\mu-a^\mu\sigma
 +{1\over2}\dot a^\mu\sigma^2
 +O(\sigma^3),
\end{align*}
更高一阶的世界线导数只出现在被略去的项中。于是
\begin{equation*}
 R^\mu(\tau')
 =\xi^\mu+u^\mu\sigma
 -{1\over2}a^\mu\sigma^2
 +{1\over6}\dot a^\mu\sigma^3
 +O(\sigma^4).
\end{equation*}
近世界线时 $\xi=O(\ell)$,而 $\sigma$ 为同阶的固有时差($c$ 的幂次只补量纲,不影响计阶)。为了把场算到有限阶,需要把统一场式中的分子、分母和光锥根一致展开。

先计算反对称分子。把世界线展开与 $R^\mu$ 的展开逐项相乘,按 $\sigma$ 幂次收集
\begin{align*}
 R^{[\mu}u^{\nu]}(\tau')={}&\xi^{[\mu}u^{\nu]}
 +\left(u^{[\mu}u^{\nu]}-\xi^{[\mu}a^{\nu]}\right)\sigma\\
 &+\left({1\over2}\xi^{[\mu}\dot a^{\nu]}
 -u^{[\mu}a^{\nu]}
 -{1\over2}a^{[\mu}u^{\nu]}\right)\sigma^2\\
 &+\left({1\over2}u^{[\mu}\dot a^{\nu]}
 +{1\over2}a^{[\mu}a^{\nu]}
 +{1\over6}\dot a^{[\mu}u^{\nu]}\right)\sigma^3
 +O(\sigma^4).
\end{align*}
利用
\begin{equation*}
 u^{[\mu}u^{\nu]}=0,
 \qquad
 a^{[\mu}a^{\nu]}=0,
 \qquad
 a^{[\mu}u^{\nu]}=-u^{[\mu}a^{\nu]},
 \qquad
 u^{[\mu}\dot a^{\nu]}=-\dot a^{[\mu}u^{\nu]},
\end{equation*}
并注意更高一阶的世界线导数项最终只进入比所需更高的阶,上式化成
\begin{equation*}
 R^{[\mu}u^{\nu]}(\tau')
 =\xi^{[\mu}u^{\nu]}
 -\xi^{[\mu}a^{\nu]}\sigma
 +\left({1\over2}\xi^{[\mu}\dot a^{\nu]}
 -{1\over2}u^{[\mu}a^{\nu]}\right)\sigma^2
 -{1\over3}\dot a^{[\mu}u^{\nu]}\sigma^3
 +O(\ell^4).
\end{equation*}
这里 $\sigma^2$ 系数中的 $-u^{[\mu}a^{\nu]}-a^{[\mu}u^{\nu]}/2$ 合并为 $-u^{[\mu}a^{\nu]}/2$,而 $\sigma^3$ 系数中的 $a^{[\mu}a^{\nu]}$ 为零、$u^{[\mu}\dot a^{\nu]}$ 与 $\dot a^{[\mu}u^{\nu]}$ 合并为 $-(2/3)\dot a^{[\mu}u^{\nu]}$。

再计算分母。把 $R^\mu(\tau')$ 与 $u^\mu(\tau')$ 逐项收缩:
\begin{align*}
 R\cdot u(\tau')={}&\xi\cdot u
 +\left(u^2-\xi\cdot a\right)\sigma
 +\left({1\over2}\xi\cdot\dot a
 -u\cdot a
 -{1\over2}a\cdot u\right)\sigma^2\\
 &+\left({1\over2}u\cdot\dot a
 +{1\over2}a^2
 +{1\over6}\dot a\cdot u\right)\sigma^3
 +O(\sigma^4).
\end{align*}
利用近点条件与运动学恒等式,并略去只影响更高阶的世界线导数项,得到
\begin{equation*}
 R\cdot u(\tau')
 =c^2\Lambda\sigma
 +{1\over2}(\xi\cdot\dot a)\sigma^2
 -{1\over6}a^2\sigma^3
 +O(\sigma^4),
 \qquad
 \Lambda\equiv1-{\xi\cdot a\over c^2}.
\end{equation*}
逐项核对:$\sigma$ 项由 $\xi\cdot u=0$ 与 $u^2=c^2$ 给出 $c^2-\xi\cdot a$;$\sigma^2$ 项中 $u\cdot a$ 的两次出现相互抵消,只剩 $\xi\cdot\dot a/2$;$\sigma^3$ 项的系数 $u\cdot\dot a/2+a^2/2+\dot a\cdot u/6$ 在 $u\cdot\dot a=-a^2$ 代入后为 $-a^2/6$,并不为零,这一项正是下面分母倒数展开中 $a^2\sigma^2/(6c^2\Lambda)$ 的来源。由于 $\xi\cdot a=O(\ell)$,有 $\Lambda=1+O(\ell)$。

为了随后对统一场式求导,需要先展开分母倒数。把上面的分母写成
\begin{equation*}
 R\cdot u=c^2\Lambda\sigma
 \left[1+{\xi\cdot\dot a\over2c^2\Lambda}\sigma
 -{a^2\over6c^2\Lambda}\sigma^2
 +O(\sigma^3)\right],
\end{equation*}
对括号作 $(1+X)^{-1}=1-X+X^2+\cdots$ 展开,
\begin{align*}
 {1\over R\cdot u}
 &={1\over c^2\Lambda\sigma}
 \left[
 1-{\xi\cdot\dot a\over2c^2\Lambda}\sigma
 +{a^2\over6c^2\Lambda}\sigma^2
 +{(\xi\cdot\dot a)^2\over4c^4\Lambda^2}\sigma^2
 +O(\ell^3)
 \right].
\end{align*}
这里 $(\xi\cdot\dot a)^2\sigma^2/c^4=O(\ell^4)$,在最终场的有限阶以前不产生贡献,故以后可省略。将化简后的分子与分母倒数相乘,先不取光锥根,得到
\begin{align*}
 {R^{[\mu}u^{\nu]}\over R\cdot u}
 ={1\over c^2\Lambda\sigma}\bigg[{}&
 \xi^{[\mu}u^{\nu]}
 -\xi^{[\mu}a^{\nu]}\sigma
 -{(\xi\cdot\dot a)\xi^{[\mu}u^{\nu]}\over2c^2\Lambda}\sigma\\
 &+{a^2\xi^{[\mu}u^{\nu]}\over6c^2\Lambda}\sigma^2
 +{1\over2}\xi^{[\mu}\dot a^{\nu]}\sigma^2
 -{1\over2}u^{[\mu}a^{\nu]}\sigma^2\\
 &+{(\xi\cdot\dot a)\xi^{[\mu}a^{\nu]}\over2c^2\Lambda}\sigma^2
 -{1\over3}\dot a^{[\mu}u^{\nu]}\sigma^3
 +O(\ell^4)
 \bigg].
\end{align*}
式中含 $\xi^{[\mu}a^{\nu]}\sigma^2$ 的交叉项在再求导并乘一个 $(R\cdot u)^{-1}$ 后只影响 $O(\ell)$,故可在所需精度下略去;同样,在方括号内部可把分母中的 $\Lambda$ 替换为 $1$——这只会丢弃 $O(\ell)$ 的修正——只有稍后与 $u^{[\mu}a^{\nu]}$ 相乘的线性 $\xi\cdot a$ 修正必须保留。于是
\begin{align*}
 {R^{[\mu}u^{\nu]}\over R\cdot u}
 ={1\over c^2\Lambda\sigma}\bigg[{}&
 \xi^{[\mu}u^{\nu]}
 -\xi^{[\mu}a^{\nu]}\sigma
 -{(\xi\cdot\dot a)\xi^{[\mu}u^{\nu]}\over2c^2}\sigma\\
 &+{a^2\xi^{[\mu}u^{\nu]}\over6c^2}\sigma^2
 +{1\over2}\xi^{[\mu}\dot a^{\nu]}\sigma^2
 -{1\over2}u^{[\mu}a^{\nu]}\sigma^2
 -{1\over3}\dot a^{[\mu}u^{\nu]}\sigma^3
 \bigg]+O(\ell^3).
\end{align*}
由于 $\tau'=\tau_0-\sigma$,有 $\d/\d\tau'=-\d/\d\sigma$。对上面的比值逐项求导,含 $\sigma^{-1}$ 的首项把幂次降低一级:
\begin{align*}
 {\d\over\d\tau'}
 \left({R^{[\mu}u^{\nu]}\over R\cdot u}\right)
 ={1\over c^2\Lambda}\bigg[{}&
 {\xi^{[\mu}u^{\nu]}\over\sigma^2}
 -{a^2\xi^{[\mu}u^{\nu]}\over6c^2}
 -{1\over2}\xi^{[\mu}\dot a^{\nu]}
 +{1\over2}u^{[\mu}a^{\nu]}
 +{2\over3}\dot a^{[\mu}u^{\nu]}\sigma
 \bigg]+O(\ell^2).
\end{align*}
现在再乘分母倒数。其中由外层展开的 $a^2\sigma^2/(6c^2\Lambda)$ 项与导数首项 $\xi^{[\mu}u^{\nu]}/\sigma^2$ 相乘产生的 $+a^2\xi^{[\mu}u^{\nu]}/(6c^2\Lambda)$,恰好抵消导数式中显式的 $-a^2\xi^{[\mu}u^{\nu]}/(6c^2)$;两者在乘以外层因子后均只差 $O(\ell)$,落在已声明的精度之内。因此在尚未代入光锥根时,场化成
\begin{align*}
 F_\eta^{\mu\nu}
 =\eta{q\over4\pi\varepsilon_0c^5\Lambda^2}
 \bigg[{}&
 {\xi^{[\mu}u^{\nu]}\over\sigma^3}
 -{(\xi\cdot\dot a)\xi^{[\mu}u^{\nu]}\over2c^2\sigma^2}
 -{\xi^{[\mu}\dot a^{\nu]}\over2\sigma}
 +{u^{[\mu}a^{\nu]}\over2\sigma}
 +{2\over3}\dot a^{[\mu}u^{\nu]}
 \bigg]_{\sigma=\sigma_\eta}
 +O(\ell).
\end{align*}

接下来必须求真正的推迟与超前根;如果简单令 $\sigma=\pm\ell/c$,会漏掉与发散幂相乘后留下的有限项。由 $R^\mu$ 的展开,
\begin{equation*}
 R^2
 =\xi^2+2\xi\cdot u\,\sigma
 +\left(u^2-\xi\cdot a\right)\sigma^2
 +\left({1\over3}\xi\cdot\dot a-u\cdot a\right)\sigma^3
 +\left({1\over4}a^2+{1\over3}u\cdot\dot a\right)\sigma^4
 +O(\ell^5).
\end{equation*}
代入 $\xi^2=-\ell^2$、$\xi\cdot u=0$、$u^2=c^2$、$u\cdot a=0$ 和 $u\cdot\dot a=-a^2$:一次项消失,$\sigma^3$ 系数中的 $u\cdot a$ 消失,而 $\sigma^4$ 系数化为 $a^2/4-a^2/3=-a^2/12$。光锥条件化成
\begin{equation*}
 0=R^2
 =-\ell^2+c^2\Lambda\sigma^2
 +{1\over3}(\xi\cdot\dot a)\sigma^3
 -{1\over12}a^2\sigma^4
 +O(\ell^5).
\end{equation*}
令 $\eta=+1$ 表示推迟根、$\eta=-1$ 表示超前根,并作 ansatz
\begin{equation*}
 \sigma_\eta
 =\eta{\ell\over c\sqrt{\Lambda}}
 \left(1+A_\eta\ell+B\ell^2+O(\ell^3)\right).
\end{equation*}
由上式,
\begin{align*}
 \sigma_\eta^2
 &={\ell^2\over c^2\Lambda}
 \left[1+2A_\eta\ell+(2B+A_\eta^2)\ell^2+O(\ell^3)\right],\\
 \sigma_\eta^3
 &=\eta{\ell^3\over c^3\Lambda^{3/2}}
 \left[1+3A_\eta\ell+O(\ell^2)\right],\\
 \sigma_\eta^4
 &={\ell^4\over c^4\Lambda^2}
 \left[1+O(\ell)\right].
\end{align*}
代入光锥条件:首项 $-\ell^2+c^2\Lambda\sigma_\eta^2$ 中的 $\ell^2$ 部分相消,除以 $\ell^2$ 后得到
\begin{align*}
 0={}&2A_\eta\ell
 +(2B+A_\eta^2)\ell^2
 +\eta{\xi\cdot\dot a\over3c^3\Lambda^{3/2}}\ell
 \left(1+3A_\eta\ell\right)
 -{a^2\over12c^4\Lambda^2}\ell^2
 +O(\ell^3).
\end{align*}
注意 $\xi\cdot\dot a=O(\ell)$,它保证组合 $(\xi\cdot\dot a)/c^3$ 的量纲与量级正确,所以 $A_\eta=O(\ell)$;因而 $A_\eta^2\ell^2$ 和 $(\xi\cdot\dot a)A_\eta\ell^2/c^3$ 都是 $O(\ell^4)$,超出当前精度。逐阶比较:由 $O(\ell)$ 阶,
\begin{equation*}
 2A_\eta+\eta{\xi\cdot\dot a\over3c^3\Lambda^{3/2}}=0,
\end{equation*}
由 $O(\ell^2)$ 阶(略去 $A_\eta^2$ 与 $(\xi\cdot\dot a)A_\eta/c^3$),
\begin{equation*}
 2B-{a^2\over12c^4\Lambda^2}=0,
\end{equation*}
所以
\begin{equation*}
 A_\eta=-\eta{\xi\cdot\dot a\over6c^3\Lambda^{3/2}},
 \qquad
 B={a^2\over24c^4\Lambda^2},
\end{equation*}
\begin{equation*}
 \sigma_\eta
 =\eta{\ell\over c\sqrt{\Lambda}}
 \left[
 1-\eta{\xi\cdot\dot a\over6c^3\Lambda^{3/2}}\ell
 +{a^2\over24c^4\Lambda^2}\ell^2
 +O(\ell^3)
 \right].
\end{equation*}
推迟根 $\sigma_{+}>0$,超前根 $\sigma_{-}<0$,与 $\tau'=\tau_0-\sigma$ 的定义一致。

场式含 $\sigma^{-3}$、$\sigma^{-2}$ 和 $\sigma^{-1}$,故把 $(1+A_\eta\ell+B\ell^2)^{-k}$ 按二项式展开到 $\ell^2$ 阶($A_\eta^2$ 部分为 $O(\ell^4)$ 可略),分别得到
\begin{align*}
 \sigma_\eta^{-1}
 &=\eta{c\sqrt{\Lambda}\over\ell}
 \left[
 1+\eta{\xi\cdot\dot a\over6c^3\Lambda^{3/2}}\ell
 -{a^2\over24c^4\Lambda^2}\ell^2
 +O(\ell^3)
 \right],\\
 \sigma_\eta^{-2}
 &={c^2\Lambda\over\ell^2}
 \left[
 1+\eta{\xi\cdot\dot a\over3c^3\Lambda^{3/2}}\ell
 -{a^2\over12c^4\Lambda^2}\ell^2
 +O(\ell^3)
 \right],\\
 \sigma_\eta^{-3}
 &=\eta{c^3\Lambda^{3/2}\over\ell^3}
 \left[
 1+\eta{\xi\cdot\dot a\over2c^3\Lambda^{3/2}}\ell
 -{a^2\over8c^4\Lambda^2}\ell^2
 +O(\ell^3)
 \right].
\end{align*}
现在逐项代入场式。首先,$\sigma^{-3}$ 项给出
\begin{align*}
 \eta{q\over4\pi\varepsilon_0c^5\Lambda^2}
 {\xi^{[\mu}u^{\nu]}\over\sigma_\eta^3}
 ={q\over4\pi\varepsilon_0c}
 \Lambda^{-1/2}\bigg[
 {\xi^{[\mu}u^{\nu]}\over c\ell^3}
 +\eta{(\xi\cdot\dot a)\xi^{[\mu}u^{\nu]}\over2c^4\Lambda^{3/2}\ell^2}
 -{a^2\xi^{[\mu}u^{\nu]}\over8c^5\Lambda^2\ell}
 \bigg]+O(\ell).
\end{align*}
其次,$\sigma^{-2}$ 项给出
\begin{equation*}
 -\eta{q\over4\pi\varepsilon_0c^5\Lambda^2}
 {(\xi\cdot\dot a)\xi^{[\mu}u^{\nu]}\over2c^2\sigma_\eta^2}
 =-{q\over4\pi\varepsilon_0c}
 \Lambda^{-1/2}
 {\eta(\xi\cdot\dot a)\xi^{[\mu}u^{\nu]}\over2c^4\Lambda^{1/2}\ell^2}
 +O(1).
\end{equation*}
两式的 $\ell^{-2}$ 项分别为 $+\eta q(\xi\cdot\dot a)\xi^{[\mu}u^{\nu]}/(8\pi\varepsilon_0c^5\Lambda^2\ell^2)$ 与 $-\eta q(\xi\cdot\dot a)\xi^{[\mu}u^{\nu]}/(8\pi\varepsilon_0c^5\Lambda\ell^2)$;将 $\Lambda^{-1}=1+\xi\cdot a/c^2+O(\ell^2)$ 一致展开后,两者在 $O(1)$ 阶恰好相消,残差 $O(\ell)$ 落在已声明的精度之内。这个抵消必须同时使用光锥根的加速度修正(它产生 $\sigma^{-3}$ 展开中的 $\ell$ 修正项)和场的显式 $\sigma^{-2}$ 项;若直接令 $\sigma_\eta=\eta\ell/c$,便看不到这一结构。

对两个 $\sigma^{-1}$ 项,使用逆幂展开并保留 $\Lambda=1-\xi\cdot a/c^2$ 的线性修正,得到
\begin{align*}
 -\eta{q\over4\pi\varepsilon_0c^5\Lambda^2}
 {\xi^{[\mu}\dot a^{\nu]}\over2\sigma_\eta}
 &=-{q\over4\pi\varepsilon_0c}
 \Lambda^{-1/2}
 {\xi^{[\mu}\dot a^{\nu]}\over2c^3\ell}
 +O(1),\\
 \eta{q\over4\pi\varepsilon_0c^5\Lambda^2}
 {u^{[\mu}a^{\nu]}\over2\sigma_\eta}
 &={q\over4\pi\varepsilon_0c}
 \Lambda^{-1/2}
 {(1+\xi\cdot a/c^2)u^{[\mu}a^{\nu]}\over2c^3\ell}
 +O(1).
\end{align*}
最后,有限项保留分支符号:
\begin{equation*}
 \eta{q\over4\pi\varepsilon_0c^5\Lambda^2}
 {2\over3}\dot a^{[\mu}u^{\nu]}
 =\eta{q\over4\pi\varepsilon_0c}
 \Lambda^{-1/2}
 {2\over3c^4}\dot a^{[\mu}u^{\nu]}
 +O(\ell).
\end{equation*}
合并以上各式,推迟场和超前场的近线展开化成
\begin{align*}
 F_\eta^{\mu\nu}
 ={q\over4\pi\varepsilon_0c}\Lambda^{-1/2}
 \bigg[{}&
 {\xi^{[\mu}u^{\nu]}\over c\ell^3}
 -{a^2\xi^{[\mu}u^{\nu]}\over8c^5\ell}
 -{\xi^{[\mu}\dot a^{\nu]}\over2c^3\ell}
 +{(1+\xi\cdot a/c^2)u^{[\mu}a^{\nu]}\over2c^3\ell}
 +\eta{2\over3c^4}\dot a^{[\mu}u^{\nu]}
 \bigg]+O(\ell).
\end{align*}
上面的展开式是前面冗长计算的简洁形式,并且仍描述同一个 $F_\eta^{\mu\nu}$。现在可以直接读出:所有发散项都与 $\eta$ 无关,只有有限的最后一项在推迟和超前两支间变号。因此
\begin{align*}
 F_{\mathrm S}^{\mu\nu}
 &={F_{+}^{\mu\nu}+F_{-}^{\mu\nu}\over2}\\
 &={q\over4\pi\varepsilon_0c}\Lambda^{-1/2}
 \left[
 {\xi^{[\mu}u^{\nu]}\over c\ell^3}
 -{a^2\xi^{[\mu}u^{\nu]}\over8c^5\ell}
 -{\xi^{[\mu}\dot a^{\nu]}\over2c^3\ell}
 +{(1+\xi\cdot a/c^2)u^{[\mu}a^{\nu]}\over2c^3\ell}
 \right]+O(\ell),
\end{align*}
而两支之差为
\begin{equation*}
 F_{+}^{\mu\nu}-F_{-}^{\mu\nu}
 ={q\over4\pi\varepsilon_0c}\Lambda^{-1/2}
 {4\over3c^4}\dot a^{[\mu}u^{\nu]}+O(\ell),
\end{equation*}
取 $\ell\to0$ 时 $\Lambda\to1$,故
\begin{equation*}
 F_{\mathrm R}^{\mu\nu}(z)
 ={1\over2}\lim_{\ell\to0}
 \left(F_{+}^{\mu\nu}-F_{-}^{\mu\nu}\right)
 ={q\over6\pi\varepsilon_0c}
 {\dot a^{[\mu}u^{\nu]}\over c^4}
 ={q\over6\pi\varepsilon_0c^5}
 \left(\dot a^\mu u^\nu-\dot a^\nu u^\mu\right).
\end{equation*}
于是同一个实际场的等价改写链现在具有明确的局部内容:$F_{\mathrm S}$ 是时间对称且携带全部 Coulomb 型奇性的部分,$F_{\mathrm{reg}}=F_{\mathrm{in}}+F_{\mathrm R}$ 在世界线上有限。奇异场携带的四动量将交给世界管上的守恒量计算。
